\section{Weak entropy and an independent sparse estimate}
\label{l:forests}

A complementary marginal may have relative entropy from uniform of
order $(\log n)^4$ and still be far from uniform in total variation.
Such information is
nevertheless sufficient for representations whose dimensions are large
enough to absorb a growing coset cap.  This section proves that assertion
and combines it with an independent Fourier estimate for representations
with a long first row.  The second law must also annihilate every
representation whose diagram has more than $n/2$ rows.  These two
Fourier properties are explicit hypotheses; the construction of the
second law plays no role in the transfer.

Throughout this section, $16$ divides $n$ and
$[n]=B_1\sqcup\cdots\sqcup B_{16}$ is a fixed partition into equal
blocks.  Let $H_j=S_{B_j}$, acting identically off $B_j$.
For a probability $\mu$ on $S_n$, let $\mu_j$ be its pushforward to
the cosets $S_n/H_j$, and let $U_j$ be uniform on this quotient.
The coset $xH_j$ specifies exactly the images under $x$ of all labels
outside $B_j$.

\subsection{Clipping a weak entropy bound}

\begin{lemma}[One retained law for sixteen entropy bounds]
\label{l:forest-clipping}
Suppose
\begin{equation}
 E_n:=\sum_{j=1}^{16}D(\mu_j\Vert U_j)
       \le C(1+\log n)^4
 \label{l:forest-entropy-hyp}
\end{equation}
for a constant $C$ independent of $n$.  Set $b_n=n^{1/50}$ and
$\eta_n=(E_n+16/e)/b_n$.  For all sufficiently large $n$, there is
an event $\mathcal G\subseteq S_n$ with
$z:=\mu(\mathcal G)\ge1-\eta_n>0$ such that
$\nu=\mu(\,\cdot\mid\mathcal G)$ satisfies
\begin{equation}
 \|\nu-\mu\|_{\mathrm{TV}}=1-z\le\eta_n=o(1),
 \qquad
 \nu(xH_j)\le\frac{M_n}{[S_n:H_j]},
 \quad M_n=2\exp(n^{1/50}).
 \label{l:forest-caps}
\end{equation}
The cap holds for every $x\in S_n$ and every $j$ under this same law
$\nu$.
\end{lemma}

\begin{proof}
Apply the clipping construction of Theorem~\ref{l:replica-clipping}
with $q=16$, $w=\mu$, $K=E_n$, and $t=b_n$. It restricts $\mu$
to the event $\mathcal G$ on which all sixteen original quotient
densities are at most $e^{b_n}$. The retained subprobability has
mass $z\ge1-\eta_n$ and satisfies every coset cap with constant
$e^{b_n}$. Here $\eta_n=o(1)$ by \eqref{l:forest-entropy-hyp},
so eventually $z\ge1/2$. Normalizing this one retained measure
increases all caps by the same factor $z^{-1}\le2$, giving
\eqref{l:forest-caps}. Since the construction is an event restriction,
its normalization is exactly $\mu(\,\cdot\mid\mathcal G)$, at
total-variation distance $1-z$ from $\mu$.
\end{proof}

\subsection{The representation range controlled by these caps}

Write $D_\lambda$ for the dimension of the irreducible representation
$V_\lambda$ of $S_n$, and put $k=n-\lambda_1$.  Here $k$ measures the
number of boxes below the first row, without transposing the diagram.
The restriction on the number of rows in the next lemma prevents a shape with a long
first column from having small dimension while $k$ is large.

\begin{lemma}[Dimension with at most $n/2$ rows]
\label{l:forest-degree}
There is an absolute $c>0$ such that, for all sufficiently large $n$,
every $\lambda\vdash n$ with $\lambda'_1\le n/2$ and $k>0$ satisfies
\begin{equation}
 \log D_\lambda\ge ck.
 \label{l:forest-degree-bound}
\end{equation}
\end{lemma}

\begin{proof}
If both the first row and the first column are shorter than $n/8$,
every hook has length less than $n/4$.  The hook formula then gives
$D_\lambda\ge(4/e)^n$, which suffices.
Otherwise transpose if necessary so the first row has length
$m=\max(\lambda_1,\lambda'_1)\ge n/8$.
The number $n-m$ of remaining boxes is at least $k/2$: it is $k$
without transposition, and after transposition the assumed bound on the number of rows
gives $n-m\ge n/2\ge k/2$.

Retain this first row and a Young subdiagram of
$s=\min(n-m,\lfloor m/2\rfloor)$ boxes below it.  For large $n$,
$s\ge k/17$.  To construct standard tableaux of the retained diagram,
put $1,\ldots,s$ in the first $s$ boxes of its top row.  Choose any
$s$ of the remaining $m$ labels for the lower diagram, placed in a fixed
standard order; fill the remaining top row in increasing order.
Every column of the lower diagram is under one of the first $s$ top
boxes, so these are valid tableaux.  There are
$\binom ms\ge(m/s)^s\ge2^s$ choices, and each extends to the full
diagram by successively adding the omitted boxes in a Young order.
Thus $D_\lambda\ge2^s$, proving the lemma.
\end{proof}

\begin{proposition}[Sixteen groups in the large-degree range]
\label{l:forest-bulk}
Let $\nu$ be any probability on $S_n$ satisfying the sixteen coset
caps in \eqref{l:forest-caps}.  Uniformly over all partitions with
\[
 k=n-\lambda_1>n^{1/20},\qquad \lambda'_1\le n/2,
\]
one has, for sufficiently large $n$,
\begin{equation}
 \|\widehat\nu(\lambda)\|_{\mathrm{op}}
 \le D_\lambda^{-1/3}.
 \label{l:forest-bulk-bound}
\end{equation}
\end{proposition}

\begin{proof}
Put $D=D_\lambda$ and restrict $V_\lambda$ to
$H_1\times\cdots\times H_{16}$.  Its irreducible types are tensor
products of sixteen carriers, with arbitrary multiplicities.
For an occurring tensor type, the product of its carrier dimensions
is at most $D$, so at least one carrier has dimension at most
$D^{1/16}$.  Assign the whole isotypic component of that type to
one such index $j$.  This decomposes every unit vector $v$ into an
orthogonal sum $v=\sum_{j=1}^{16}v_j$, where $v_j$ is supported only
on $H_j$-types of dimension at most $D^{1/16}$.

Branching shows that each $H_j$-type diagram is a subdiagram of
$\lambda$; its tail below its first row has at most $k$ boxes.
Let $p(i)$ count the partitions of $i$, with $p(0)=1$.  There are at most
\[
 B_k=\sum_{i=0}^k p(i),\qquad
 \log B_k=O(\sqrt k\log(k+1)),
\]
such types, by the elementary multiplicity-and-list count in
\eqref{l:partition-counts}.  Its sharper exponential bound also
suffices.

Let $W_j$ be the span of the $H_j$-orbit of $v_j$.
For one type of degree $a$, collect all equivalent copies as
$V_a\otimes\mathcal M_a$.  The orbit of the component of $v_j$
lies in the image of
\[
 \operatorname{End}(V_a)\longrightarrow V_a\otimes\mathcal M_a,
 \qquad T\longmapsto(T\otimes I)v_j.
\]
Its dimension is at most $a^2$, regardless of the multiplicity.
Consequently
\begin{equation}
 \dim W_j\le B_k D^{1/8}.
 \label{l:forest-orbit-span}
\end{equation}
This is the point at which grouping equivalent copies is essential:
counting each copy separately would lose the degree saving.

We now use the caps to average these small orbit spaces over the full
group.  Since $W_j$ is $H_j$-invariant, the space
$\rho_\lambda(x)W_j$ depends only on $xH_j$.
Its projection, averaged uniformly over the group, is
$(\dim W_j/D)I$ by Schur's lemma and the trace.
For a unit vector $w$, Cauchy--Schwarz and the coset cap therefore give
\[
 \begin{split}
 |\langle w,\widehat\nu(\lambda)v_j\rangle|
 &\le\|v_j\|
 \left(\mathbb E_{x\sim\nu}
       \|P_{\rho_\lambda(x)W_j}w\|^2\right)^{1/2}\\
 &\le\|v_j\|\sqrt{M_n B_kD^{-7/8}}.
 \end{split}
\]
Because $\sum_j\|v_j\|\le4$, summing yields the quantitative bound
\begin{equation}
 \|\widehat\nu(\lambda)\|_{\mathrm{op}}
 \le4\sqrt{M_n B_k}\,D^{-7/16}.
 \label{l:forest-unabsorbed}
\end{equation}

It remains to absorb the growing cap and the type count.
Lemma~\ref{l:forest-degree} gives $\log D\ge ck$, whereas
$\log M_n=O(n^{1/50})$ and
$\log B_k=O(\sqrt k\log(k+1))$.
For $k>n^{1/20}$, both of the latter quantities are $o(k)$
uniformly.  Thus, for large $n$,
$\log(16M_nB_k)\le(5/24)\log D$.
Squaring \eqref{l:forest-unabsorbed} and using this inequality gives
$16M_nB_kD^{-7/8}\le D^{-2/3}$, as required.
\end{proof}

\subsection{Combining the large and small degrees}

The preceding proposition leaves two classes untreated: diagrams with
more than $n/2$ rows, and diagrams with at most $n^{1/20}$ boxes below
the first row.  The following theorem states exactly what a second law
must provide on those classes.

\begin{theorem}[Entropy and sparse-degree hybrid]
\label{l:forest-hybrid}
Let $n\to\infty$ through multiples of $16$.
For each $n$, let $\mu$ and $K$ be probabilities on $S_n$.
Suppose a fixed equal sixteen-block partition satisfies
\eqref{l:forest-entropy-hyp} for $\mu$.  Suppose also that $K$ satisfies
\begin{align}
 \|\widehat K(\lambda)\|_{\mathrm{op}}
 &\le n^{-k/10}
 &&\left(1\le k=n-\lambda_1\le n^{1/20}\right),
 \label{l:forest-sparse-hyp}\\
 \widehat K(\lambda)&=0
 &&\left(\lambda'_1>n/2\right).
 \label{l:forest-row-hyp}
\end{align}
Then the retained probability $\nu$ from
Lemma~\ref{l:forest-clipping} satisfies
\begin{equation}
 \|\nu^{*12}*K^{*40}-U_{S_n}\|_{\mathrm{TV}}=o(1),
 \qquad
 \|\mu^{*12}*K^{*40}-U_{S_n}\|_{\mathrm{TV}}=o(1).
 \label{l:forest-hybrid-conclusion}
\end{equation}
All factors in these convolutions are independent.  There is no assumed
identity between $K$ and a convolution root or power of $\mu$.
\end{theorem}

\begin{proof}
Put $\sigma=\nu^{*12}*K^{*40}$.  The Fourier matrices multiply in
this order.  Each factor is a contraction, since it is an average of
unitary matrices.  For $\lambda'_1\le n/2$ and
$k>n^{1/20}$, Proposition~\ref{l:forest-bulk} gives
$\|\widehat\sigma(\lambda)\|_{\mathrm{op}}\le D_\lambda^{-4}$.
For more than $n/2$ rows, \eqref{l:forest-row-hyp} gives zero.
For $1\le k\le n^{1/20}$, \eqref{l:forest-sparse-hyp} gives
$\|\widehat\sigma(\lambda)\|_{\mathrm{op}}\le n^{-4k}$.
These statements use submultiplicativity and do not require normality
of any Fourier matrix.

Plancherel bounds four times the squared total-variation distance by
$\sum_{\lambda\ne(n)}D_\lambda^2
\|\widehat\sigma(\lambda)\|_{\mathrm{op}}^2$.
For each fixed $k$, there are at most $p(k)$ diagrams with
$n-\lambda_1=k$, since the boxes below the first row form a
partition of $k$.  By Lemma~\ref{l:forest-degree}, the contribution
from the large-degree range is at most
\[
 \sum_{k>n^{1/20}}p(k)e^{-6ck}=o(1).
\]
For the small-degree range, branching shows that $V_\lambda$ occurs
in the permutation representation on ordered distinct $k$-tuples:
its restriction to $S_{n-k}$ contains the trivial representation
because its first row has length $n-k$.
Thus $D_\lambda\le(n)_k\le n^k$, and that contribution is at most
\[
 \sum_{1\le k\le n^{1/20}}p(k)n^{2k}n^{-8k}
 \le\sum_{k\ge1}e^{3\sqrt k}n^{-6k}=o(1).
\]
This also explains why no separate sign hypothesis is needed here:
the sign diagram has more than $n/2$ rows and is annihilated by $K$.
The first assertion follows.  Convolution contracts total variation,
so replacing twelve factors changes the distance by at most
$12\|\mu-\nu\|_{\mathrm{TV}}=o(1)$, proving the second.
\end{proof}
